\begin{document}

\section{Analysis of the problem}
We decompose the object into five classes; the classification is an anatomy of
the object and we ask where does each number come from. The structural,
configuration, temporal, tabulated and computed classes differ in what a checker
can ask, which is why the decomposition decides where the checks go. A checker
that treats every number alike measures the layout rather than the provenance,
and the failure rate it reports is therefore mostly noise about formatting: in a
sample of 23 papers the layout share reached 9.4 percent of all tokens.

\begin{table}\centering
\begin{tabular}{@{}lll@{}}
\toprule
Class & Source & Share (\%) \\
\midrule
structural & layout & 9.4 \\
configuration & method & 0.0 \\
temporal & input file & 14.0 \\
computed & artifact & 71.1 \\
\bottomrule
\end{tabular}
\end{table}

\section{Method and why}
Each mechanism carries a rejected alternative; we explain why this rather than
that, and what would change our mind before we kept it. The alternative we
rejected for each is stated in the table, together with the observation that
would have overturned the choice, and the measured cost of keeping it.

\begin{table}\centering
\begin{tabular}{@{}llll@{}}
\toprule
Mechanism & Rejected alternative & Why & Cost (s) \\
\midrule
freeze & inline quotes & checkable by a third party & 0.4 \\
classify & trace every token & inflates with layout & 1.2 \\
veto & a single score & five channels outvote two & 2.6 \\
\bottomrule
\end{tabular}
\end{table}

\section{What was wrong}
\subsection{Case 1: the tokenizer}
It was wrong; how found: a count of 199 against 1293; the change; why this is
right, because the count moved while the reports did not.
\subsection{Case 2: the gate}
Wrong again; how found: a paper with zero citations passed; the change; why
this is right, because a check must be able to fail on its own object.

\section{Deployment}
Three modes, a playbook, a time budget and a cost that separates machine time
from editorial time. The first mode is a command line run by the author on the
competition laptop, the second is a service any agent can call, and the third is
a gate that runs on every commit. The failure table is the part that actually
gets used under deadline, because the failure that matters is not a wrong
answer but a stalled pipeline with a report due in an hour; a stalled compile
costs 12 minutes of editorial time and nothing in machine time.

\begin{table}\centering
\begin{tabular}{@{}lll@{}}
\toprule
Mode & Machine (min) & Editorial (min) \\
\midrule
cli & 0.1 & 0.0 \\
service & 0.3 & 2.0 \\
ci & 1.5 & 0.0 \\
\bottomrule
\end{tabular}
\end{table}

\end{document}
